\chapter{Đăng ký thi}

Chương này gồm \textbf{7 đường dẫn} của mô-đun \rt{wujia_portal_exam}, gom thành \textbf{7 mục}
(cộng một mục cho đường dẫn cũ). Một đường dẫn --- \rt{/portal/exam/register} --- \textbf{có hai
hàm}: một hàm mở màn, một hàm nhận phiếu gửi lên; chương này tách thành hai mục vì hai việc khác
hẳn nhau.

\textbf{Bốn điểm cần biết trước khi đọc:}

\begin{enumerate}[leftmargin=*]
  \item \textbf{Trên UAT hiện KHÔNG ai đăng ký thi được.} Có \uat{3} khóa thi đã phát hành và
        \uat{2} kỳ thi còn ở trạng thái mở, nhưng ngày thi của chúng là \textbf{16/08} và
        \textbf{25/08} --- đã qua --- và hạn đăng ký là 09:00 chính ngày thi. Vì vậy cả ba khóa
        đều hiện chữ \emph{``Đã đóng''} và lịch \textbf{không có ngày nào bấm được}. Giống hệt
        tình trạng của Đổi trả ở chương 6, cần dựng dữ liệu mẫu trước khi nghiệm thu.
  \item \textbf{Không chặn vai trò.} Nhân viên đăng ký thi được như Chủ tiệm. Đây là phân hệ
        ghi dữ liệu duy nhất ngoài Đặt hàng không có cổng vai trò.
  \item \textbf{Phạm vi là CỬA HÀNG ĐANG CHỌN, không phải mọi cửa hàng.} Người phụ trách nhiều
        cửa hàng phải bấm đổi cửa hàng mới xem được phiếu của cửa hàng kia --- giống Công nợ,
        khác Đổi trả và Yêu cầu thông tin.
  \item \textbf{Gửi phiếu xong là chốt.} Portal không có nút sửa, không có nút huỷ, không gửi
        lại được phiếu bị từ chối. Mọi thay đổi phải nhờ Ngô Gia làm trong ERP.
\end{enumerate}

Số thật trên UAT: \uat{1} phiếu đăng ký (của cửa hàng \emph{TP HCM Quận 1}, do Administrator
tạo, đã được xác nhận, \uat{1} người dự thi, kết quả \emph{Đạt}), \uat{3} kỳ thi, \uat{3} ca giờ,
\uat{0} phiếu chờ duyệt, \uat{0} phiếu bị từ chối, \uat{0} phiếu bị huỷ.

% =============================================================================
\section{Màn Danh sách phiếu đăng ký thi}
\rt{/portal/exam}
\medskip

\mvao

Mục \emph{Đăng ký thi} trên thanh điều hướng. Tham số: \rt{page}, \rt{limit}, \rt{state},
\rt{result}, \rt{q}, \rt{date_from}, \rt{date_to}.

\mai

\begin{itemize}[leftmargin=*]
  \item Phải đăng nhập. \textbf{Không chặn vai trò} --- mọi thành viên cửa hàng đều xem được.
  \item Chỉ thấy phiếu của \textbf{cửa hàng đang chọn}. \textbf{Chưa chọn cửa hàng} thì màn
        hiện đúng câu \emph{``Chưa có đăng ký thi''} --- \textbf{không} báo là do chưa chọn
        cửa hàng, nên người dùng dễ tưởng cửa hàng mình chưa từng đăng ký.
  \item Thực tế UAT: \uat{1} phiếu duy nhất, thuộc cửa hàng \emph{TP HCM Quận 1}. Ba cửa hàng
        còn lại mở ra màn rỗng.
\end{itemize}

\mhien

\begin{hienbang}
Danh sách phiếu &
  \rt{wujia.exam.registration} &
  Phiếu có cửa hàng \textbf{đúng bằng cửa hàng đang chọn} &
  Ngày nộp phiếu mới nhất trước; \textbf{10} phiếu mỗi trang (10/20/50) &
  \uat{1} phiếu toàn hệ \\ \addlinespace
Huy hiệu trạng thái &
  \rt{state} &
  Bốn nhãn: \emph{Chờ xác nhận} · \emph{Đã đăng ký} · \emph{Từ chối} · \emph{Đã hủy}.
  Bản điện thoại gọi nhãn đầu là \emph{``Chờ duyệt''} --- \textbf{cùng một trạng thái, hai chữ
  khác nhau trên hai giao diện}. &
  --- &
  \uat{1} Đã đăng ký · \uat{0} ở ba trạng thái còn lại \\ \addlinespace
Cột kết quả &
  \rt{wujia.exam.session} &
  Chỉ hiện khi \textbf{kỳ thi đã công bố kết quả}; khi đó đếm \emph{n Đạt · n Không đạt} trong
  phiếu. Chưa công bố thì ghi \emph{``Chưa công bố''}. &
  --- &
  \uat{1} kỳ đã công bố · \uat{1} người Đạt · \uat{0} Không đạt \\ \addlinespace
Dòng trên bản điện thoại &
  --- &
  Khi đã công bố, huy hiệu đổi thành \emph{``Có kết quả''} và \textbf{đè lên} trạng thái phiếu
  --- phiếu bị từ chối thuộc kỳ đã công bố cũng hiện \emph{``Có kết quả''} &
  --- & --- \\
\end{hienbang}

\mloc

\begin{itemize}[leftmargin=*]
  \item \textbf{Từ khoá} --- khớp một phần theo \textbf{mã phiếu} hoặc \textbf{tên khóa thi}.
        Không tìm theo tên thí sinh.
  \item \textbf{Trạng thái} --- một trong bốn nhãn; giá trị lạ bị bỏ qua.
  \item \textbf{Kết quả} --- \emph{Đã công bố} / \emph{Chưa công bố}, lọc theo cờ của
        \textbf{kỳ thi} chứ không theo từng người.
  \item \textbf{Khoảng ngày} --- lọc theo \textbf{ngày nộp phiếu}, quy đổi đúng múi giờ tài
        khoản. Ngày ngược thì \textbf{chặn truy vấn}, giữ nguyên hai ô đã gõ và báo tại thanh
        lọc --- giống Thông báo và Giao hàng, khác Công nợ (bên đó lặng lẽ đảo hai ngày).
  \item \textbf{Số dòng mỗi trang} --- 10, 20, 50; gõ số ngoài danh sách thì tự về 10.
\end{itemize}

\mkiem{Đăng ký mới}

Không kiểm gì ở màn danh sách --- mọi phép kiểm nằm ở màn đăng ký.

\mghi

\textbf{Không ghi gì.}

\mhoi

\begin{ask}
\begin{enumerate}[leftmargin=*]
  \item \textbf{Chưa chọn cửa hàng nhưng màn báo ``Chưa có đăng ký thi''} --- sai nguyên nhân.
        Nên đổi thành lời nhắc chọn cửa hàng như các màn khác.
  \item \textbf{Cùng một trạng thái, hai chữ khác nhau giữa máy tính và điện thoại}
        (\emph{Chờ xác nhận} · \emph{Chờ duyệt}). Anh chọn một chữ để thống nhất.
  \item \textbf{Huy hiệu ``Có kết quả'' đè lên trạng thái phiếu} trên bản điện thoại --- phiếu
        bị từ chối hay bị huỷ nhưng thuộc kỳ đã công bố vẫn hiện \emph{Có kết quả}. Cần sửa.
  \item \textbf{Chỉ xem được cửa hàng đang chọn}, trong khi luật dữ liệu của hệ thống cho phép
        xem mọi cửa hàng của mình. Anh muốn gộp như Đổi trả không?
\end{enumerate}
\end{ask}

% =============================================================================
\section{Màn Đăng ký thi --- chọn khóa và lịch}
\rt{/portal/exam/register} \note{(mở màn)}
\medskip

\mvao

Nút \emph{Đăng ký mới} ở màn danh sách. Tham số: \rt{course_id}, \rt{year}, \rt{month}.

\mai

\begin{itemize}[leftmargin=*]
  \item Phải đăng nhập \textbf{và phải đang chọn một cửa hàng} --- chưa chọn thì bị đưa thẳng
        về danh sách, \textbf{không có thông báo}.
  \item Không chặn vai trò.
\end{itemize}

\mhien

\begin{hienbang}
Danh sách khóa thi &
  \rt{wujia.exam.course} &
  Khóa \textbf{đã phát hành} và còn hiệu lực. Khóa còn nháp thì portal không thấy. &
  Theo tên &
  \uat{3}/3 khóa đã phát hành \\ \addlinespace
Nhãn trạng thái khóa &
  \rt{wujia.exam.session} &
  \emph{``Còn lịch''} khi còn ít nhất một kỳ đăng ký được; \emph{``Hết chỗ''} khi còn kỳ trong
  hạn nhưng hết ghế; \emph{``Đã đóng''} khi không còn kỳ nào &
  --- &
  Cả \uat{3} khóa đang hiện \textbf{Đã đóng} \\ \addlinespace
Dòng mô tả khóa &
  \rt{registration_horizon_days} &
  \emph{``n kỳ thi $\bullet$ Trong 60 ngày tới''} --- chỉ đếm kỳ \textbf{đang mở} có ngày thi
  \textbf{từ hôm nay} tới hết cửa sổ &
  --- &
  Cả 3 khóa đặt cửa sổ \uat{60} ngày; hiện đếm ra \uat{0} kỳ \\ \addlinespace
Lịch tháng &
  \rt{wujia.exam.session} &
  Kỳ thi của khóa đang chọn, trong tháng đang xem, \textbf{chưa bị huỷ}. Mỗi ngày mang một
  trong ba trạng thái: \emph{còn chỗ} · \emph{hết chỗ} · \emph{không có lịch}. &
  Tuần bắt đầu từ Thứ 2 &
  \uat{0} ngày bấm được \\ \addlinespace
Bảng tóm tắt (máy tính) &
  --- &
  Tên khóa, tên cửa hàng đang chọn, \emph{``0 / n''} số người đã thêm trên giới hạn mỗi phiếu.
  Ngày thi · địa điểm · hạn đăng ký để trống, do chọn khung giờ mới có. &
  --- & --- \\
\end{hienbang}

Một ngày được coi là \textbf{còn chỗ} khi kỳ thi hôm đó \textbf{đang mở}, \textbf{chưa quá hạn
đăng ký}, \textbf{còn ghế trống}, và ngày đó nằm trong khoảng \textbf{từ hôm nay tới hết cửa sổ
60 ngày}. Thiếu một trong bốn điều kiện thì ngày đó không bấm được.

\mloc

Chọn \textbf{khóa thi} và bấm mũi tên đổi \textbf{tháng}. Chọn khóa không có trong danh sách
đã phát hành thì màn \textbf{tự quay về khóa đầu tiên}. Gõ tháng lạ trên thanh địa chỉ có thể
làm màn báo lỗi kỹ thuật --- số tháng chưa được kiểm.

\mkiem{một ngày trên lịch}

Bấm ngày chỉ mở bảng khung giờ; mọi phép kiểm thật nằm ở bước gửi phiếu.

\mghi

\textbf{Không ghi gì.}

\mhoi

\begin{ask}
\begin{enumerate}[leftmargin=*]
  \item \textbf{UAT không có kỳ thi nào đăng ký được} --- cần tạo ít nhất một kỳ ngày tương lai
        để nghiệm thu được toàn bộ luồng.
  \item \textbf{Chưa chọn cửa hàng thì bị đá về danh sách không lời giải thích.}
  \item \textbf{Số tháng gõ trên thanh địa chỉ chưa được kiểm} --- dev sẽ đưa vào đợt chuẩn hoá,
        không cần task riêng.
  \item \textbf{``Cửa sổ 60 ngày'' đang đặt ở khóa thi}, mỗi khóa một giá trị. Anh xác nhận để
        Ngô Gia tự chỉnh theo từng khóa, hay muốn chốt chung một con số?
\end{enumerate}
\end{ask}

% =============================================================================
\section{Lấy lịch tháng của một khóa thi}
\rt{/portal/exam/calendar}
\medskip

\mvao

Không phải màn --- gọi ngầm khi người dùng đổi khóa thi hoặc bấm mũi tên sang tháng khác.

\mai

Phải đăng nhập \textbf{và đang chọn cửa hàng}; chưa chọn thì trả về lời nhắc \emph{``Vui lòng
chọn cửa hàng trước khi thao tác''}. Không chặn vai trò.

\mhien

Trả về ma trận tuần của tháng, mỗi ô một ngày kèm trạng thái, đúng như mô tả ở màn trên.

\mloc

Ba tham số: khóa thi, năm, tháng. Thiếu năm/tháng thì lấy tháng hiện tại.

\mkiem{gọi}

\begin{kiembang}
1 & Đang chọn cửa hàng & ``Vui lòng chọn cửa hàng trước khi thao tác.'' \\
2 & Khóa thi tồn tại \textbf{và đang ở trạng thái đã phát hành} & ``Khóa thi không hợp lệ.'' \\
\end{kiembang}

\textbf{Không} kiểm cửa hàng có liên quan gì tới khóa thi --- lịch thi là dữ liệu chung, cửa
hàng nào cũng xem được như nhau.

\mghi

\textbf{Không ghi gì.}

\mhoi

\begin{ask}
Không có điểm cần chốt.
\end{ask}

% =============================================================================
\section{Lấy khung giờ thi của một ngày}
\rt{/portal/exam/slots}
\medskip

\mvao

Không phải màn --- gọi ngầm khi người dùng bấm một ngày còn chỗ trên lịch.

\mai

Giống mục trên: phải đăng nhập và đang chọn cửa hàng.

\mhien

\begin{hienbang}
Danh sách khung giờ &
  \rt{wujia.exam.session} &
  Kỳ thi của khóa đang chọn, \textbf{đúng ngày đó}, chưa bị huỷ &
  Theo giờ bắt đầu &
  \uat{3} ca giờ đang bật: 08--10, 14--16, 16--17 \\ \addlinespace
Nhãn tình trạng &
  --- &
  Bốn khả năng theo thứ tự xét: \emph{Đã đóng} (kỳ không còn mở) $\rightarrow$ \emph{Hết hạn}
  (quá hạn đăng ký) $\rightarrow$ \emph{Hết chỗ} (không còn ghế) $\rightarrow$ \emph{Còn n chỗ} &
  --- & --- \\ \addlinespace
Thông tin kèm theo &
  \rt{wujia.exam.session} &
  Địa điểm, hạn đăng ký, số ghế còn, \textbf{giới hạn người mỗi phiếu} &
  --- &
  Kỳ 25/08 cho \uat{3} người/phiếu trong khi khóa của nó đặt \uat{2} \\
\end{hienbang}

\textbf{Giới hạn người mỗi phiếu} lấy của \textbf{kỳ thi} trước; kỳ để trống mới lấy của khóa.
Đây là nguồn duy nhất dùng cho cả câu hướng dẫn trên màn, phép đếm hạn mức, lẫn phép kiểm lúc
gửi --- nên ba chỗ không bao giờ lệch nhau.

\mloc

Hai tham số: khóa thi và ngày. Ngày nhận cả dạng \emph{2026-08-25} lẫn \emph{25/08/2026}.

\mkiem{gọi}

\begin{kiembang}
1 & Đang chọn cửa hàng & ``Vui lòng chọn cửa hàng trước khi thao tác.'' \\
2 & Khóa thi tồn tại và đã phát hành, ngày đọc được & ``Lịch thi không hợp lệ.'' \\
\end{kiembang}

\textbf{Không} kiểm ngày có nằm trong cửa sổ 60 ngày hay không --- màn lịch đã chặn từ trước
bằng cách không cho bấm những ngày đó.

\mghi

\textbf{Không ghi gì.}

\mhoi

\begin{ask}
\textbf{Kỳ thi đặt giới hạn người riêng, khác giới hạn của khóa.} Trên UAT kỳ 25/08 cho 3 người
mỗi phiếu còn khóa của nó ghi 2. Anh xác nhận cách này (kỳ đè khóa) là đúng ý.
\end{ask}

% =============================================================================
\section{Gửi phiếu đăng ký thi}
\rt{/portal/exam/register} \note{(gửi phiếu)}
\medskip

\mvao

Không phải màn --- nút \emph{Xác nhận đăng ký} ở bước cuối của màn đăng ký.

\mai

Phải đăng nhập và đang chọn cửa hàng. Không chặn vai trò. Cửa hàng ghi vào phiếu là
\textbf{cửa hàng đang chọn}, không phải số mà màn gửi lên --- nên không sửa được bằng cách can
thiệp dữ liệu gửi đi.

\mhien

Gửi xong thì chuyển thẳng sang màn chi tiết của phiếu vừa tạo.

\mloc

\begin{itemize}[leftmargin=*]
  \item \textbf{Khung giờ} --- chọn từ bảng khung giờ.
  \item \textbf{Danh sách người dự thi} --- mỗi người gồm \textbf{họ tên} (bắt buộc),
        \textbf{số điện thoại} (bắt buộc), \textbf{năm sinh}, \textbf{chức vụ}, \textbf{ảnh}.
        Người dự thi \textbf{gõ tay hoàn toàn}, không chọn từ danh sách nhân sự nào.
  \item \textbf{Ghi chú} --- một dòng, không bắt buộc.
  \item \textbf{Ảnh} --- chỉ nhận PNG hoặc JPEG, tối đa \textbf{5 MB} mỗi ảnh; ảnh được
        \textbf{thu nhỏ về tối đa 1920 điểm ảnh} trước khi lưu.
\end{itemize}

\mkiem{Xác nhận đăng ký}

\begin{kiembang}
1 & Đang chọn cửa hàng & ``Vui lòng chọn cửa hàng trước khi thao tác.'' \\
2 & Khung giờ còn tồn tại và khóa của nó đã phát hành &
    ``Khung giờ đã thay đổi hoặc không còn. Vui lòng chọn lại lịch thi.'' \\
3 & Có ít nhất 1 người dự thi & ``Cần ít nhất 1 người dự thi.'' \\
4 & Số người không vượt giới hạn mỗi phiếu & ``Tối đa n người mỗi phiếu.'' \\
5 & Mỗi người đủ họ tên và số điện thoại & ``Mỗi người dự thi cần có họ tên và số điện thoại.'' \\
6 & Số điện thoại đúng dạng (bắt đầu 0 hoặc +84, 9--11 chữ số) &
    ``Số điện thoại '…' không hợp lệ (vd 0901234567).'' \\
7 & Năm sinh là số & ``Năm sinh '…' không hợp lệ.'' \\
8 & Ảnh đúng định dạng, đúng dung lượng, không hỏng &
    ``Ảnh nhân viên không đúng định dạng hoặc vượt dung lượng cho phép.'' /
    ``Ảnh nhân viên không hợp lệ hoặc bị hỏng.'' \\
9 & Kỳ thi \textbf{đang mở đăng ký} & ``Kỳ thi '…' không mở đăng ký.'' \\
10 & \textbf{Chưa quá hạn đăng ký} & ``Kỳ thi '…' đã quá hạn đăng ký.'' \\
11 & \textbf{Còn đủ ghế} sau khi cộng số người của phiếu này &
    ``Kỳ thi '…' đã hết chỗ (sức chứa n, đã giữ m).'' \\
12 & Không có phiếu khác đang được ghi cùng lúc cho kỳ này &
    ``Hệ thống đang xử lý đăng ký khác cho kỳ thi này, vui lòng thử lại.'' \\
\end{kiembang}

Bước 9--12 là phần đáng chú ý nhất: \textbf{hệ thống khoá dòng kỳ thi lại rồi mới đếm ghế},
nên hai cửa hàng bấm gửi cùng lúc không thể cùng chiếm chỗ cuối. Nếu kỳ thi hết chỗ ngay lúc
đó, phiếu \textbf{không được ghi} và người dùng nhận câu báo rõ ràng thay vì lỗi kỹ thuật.

Năm sinh còn bị kiểm thêm ở tầng dữ liệu: phải nằm trong khoảng \textbf{1900 đến năm hiện tại}.

\mghi

\begin{itemize}[leftmargin=*]
  \item Tạo một phiếu \rt{wujia.exam.registration}, mã tự sinh dạng \emph{WJ-EXR/26/00004},
        trạng thái \textbf{Chờ xác nhận} --- \textbf{không có trạng thái nháp}.
  \item Người đăng ký = người đang đăng nhập; cửa hàng = cửa hàng đang chọn; hệ thống còn
        \textbf{tự tìm tư cách thành viên} của người đó tại cửa hàng để gắn vào phiếu.
  \item Tạo một dòng \rt{wujia.exam.registration.line} cho mỗi người dự thi, kết quả để
        \emph{Chưa có}.
  \item \textbf{Ghi ngay số ghế đã giữ} của kỳ thi --- phiếu chờ duyệt đã chiếm chỗ, không đợi
        Ngô Gia xác nhận.
  \item \textbf{Không} gửi email, \textbf{không} tạo việc cần làm. Ngô Gia phải tự vào ERP xem.
\end{itemize}

\mhoi

\begin{ask}
\begin{enumerate}[leftmargin=*]
  \item \textbf{Phiếu chờ duyệt đã chiếm ghế.} Cửa hàng gửi phiếu rồi bỏ đó sẽ khoá chỗ của
        cửa hàng khác cho tới khi Ngô Gia từ chối. Anh muốn giữ vậy, hay chỉ tính ghế khi đã
        xác nhận?
  \item \textbf{Người dự thi gõ tay hoàn toàn}, không nối với danh sách nhân sự của cửa hàng.
        Gõ sai tên thì không có gì đối chiếu. Anh có cần nối vào danh sách thành viên cửa hàng
        không?
  \item \textbf{Gửi phiếu xong không ai được báo} --- giống Hỗ trợ ở chương 7.
  \item \textbf{Không có bước xem lại trước khi gửi ở bản điện thoại} sau khi bấm xác nhận ---
        bấm là gửi luôn.
\end{enumerate}
\end{ask}

% =============================================================================
\section{Màn Chi tiết phiếu đăng ký · kết quả thi}
\rt{/portal/exam/registration/<int:reg_id>}
\medskip

\mvao

Bấm một dòng ở màn danh sách, hoặc được đưa tới ngay sau khi gửi phiếu.

\mai

Phiếu phải thuộc \textbf{cửa hàng đang chọn}. Phiếu của cửa hàng khác --- kể cả cửa hàng mà tôi
cũng truy cập được --- \textbf{không mở được} cho tới khi đổi sang cửa hàng đó. Sai thì về danh
sách, không có thông báo. Không chặn vai trò.

\mhien

\begin{hienbang}
Thông tin phiếu &
  \rt{wujia.exam.registration} &
  Mã phiếu, khóa thi, ngày giờ thi, địa điểm, cửa hàng, người đăng ký, ngày nộp, số người &
  --- &
  Phiếu UAT: \uat{1} người, nộp 03/08 \\ \addlinespace
Băng thông báo theo trạng thái &
  \rt{state} &
  Bốn trạng thái ra bốn băng khác nhau: chờ xác nhận (nhắc portal chưa hiện kết quả) · đã xác
  nhận · từ chối (kèm \textbf{lý do từ chối}) · đã huỷ (kèm \textbf{lý do huỷ}) &
  --- & --- \\ \addlinespace
Danh sách người dự thi &
  \rt{wujia.exam.registration.line} &
  Họ tên, điện thoại, năm sinh, chức vụ, ảnh &
  Theo thứ tự nhập &
  \uat{0}/1 người có ảnh \\ \addlinespace
Cột kết quả &
  \rt{result} &
  \textbf{Chỉ hiện khi kỳ thi đã công bố kết quả}; chưa công bố thì mọi người đều hiện
  \emph{Chưa có} dù Ngô Gia đã nhập điểm trong ERP &
  --- &
  \uat{1} người Đạt \\ \addlinespace
Bảng kết quả (máy tính) &
  --- &
  Chỉ hiện khi phiếu \textbf{đã được xác nhận}. Phiếu bị từ chối hoặc bị huỷ không có bảng này. &
  --- & --- \\ \addlinespace
Lời nhắc thi lại &
  --- &
  Chỉ hiện khi đã công bố \textbf{và} thật sự có người không đạt &
  --- & --- \\
\end{hienbang}

\mloc

Không có ô lọc, không có ô nhập.

\mkiem{--- }

\textbf{Màn này không có nút nào ghi dữ liệu.} Không sửa, không huỷ, không gửi lại.

\mghi

\textbf{Không ghi gì.}

\mhoi

\begin{ask}
\begin{enumerate}[leftmargin=*]
  \item \textbf{Portal không huỷ được phiếu.} Cửa hàng đăng ký nhầm phải gọi Ngô Gia; trong
        ERP người huỷ còn bắt buộc nhập lý do. Anh có muốn cho cửa hàng tự huỷ khi phiếu còn
        \emph{Chờ xác nhận} không?
  \item \textbf{Phiếu bị từ chối là ngõ cụt} --- không sửa, không gửi lại; phải tạo phiếu mới
        từ đầu, gõ lại toàn bộ người dự thi.
  \item \textbf{Chỉ mở được phiếu của cửa hàng đang chọn}, dù luật dữ liệu cho phép xem mọi cửa
        hàng của mình --- cùng câu hỏi với màn danh sách.
\end{enumerate}
\end{ask}

% =============================================================================
\section{Ảnh thí sinh trong phiếu}
\rt{/portal/exam/line/<int:line_id>/photo}
\medskip

\mvao

Không phải màn --- là ảnh hiện trong bảng người dự thi ở màn chi tiết.

\mai

\begin{kiembang}
1 & Đang chọn cửa hàng & Từ chối truy cập \\
2 & Dòng người dự thi thuộc \textbf{cửa hàng đang chọn} & Báo không tìm thấy \\
3 & Người đó thật sự có ảnh & Báo không tìm thấy \\
\end{kiembang}

\mhien

Trả về ảnh đã thu nhỏ, kèm cơ chế đệm của trình duyệt nên mở lại không tải lại.

\mloc

Không có.

\mkiem{ảnh}

Xem bảng trên.

\mghi

\textbf{Không ghi gì.}

\mhoi

\begin{ask}
Không có điểm cần chốt.
\end{ask}

% =============================================================================
\section{Đường dẫn cũ của phân hệ Đăng ký thi}
\rt{/portal/exam-registration}
\medskip

Đường dẫn của bản portal cũ (v14) được giữ lại và \textbf{chuyển hướng vĩnh viễn} sang
\rt{/portal/exam}. Lý do giữ: khoảng 1.500 tài khoản cửa hàng có thể còn lưu liên kết cũ trong
trình duyệt hoặc trong email Ngô Gia đã gửi.

Đường dẫn này \textbf{không đòi đăng nhập} --- nó chỉ chuyển hướng; màn đích mới hỏi đăng nhập.
\textbf{Không ghi gì.}

\begin{ask}
Không có điểm cần chốt --- đây là phần giữ tương thích, BA không cần lên task.
\end{ask}
